\section{Deep Reinforcement Learning as Learned Representation Plus Control}

\subsection{Overview}

Classical reinforcement learning assumes that value functions and policies can be represented exactly.  
In modern settings, state and action spaces are large or continuous, and learning must rely on \emph{function approximation}.

\medskip

\noindent
\textbf{Key idea.}  
Deep reinforcement learning (deep RL) replaces tabular representations with parameterized functions:
\[
V(s) \approx V_\theta(s), \quad Q(s,a) \approx Q_\theta(s,a), \quad \pi(a \mid s) \approx \pi_\theta(a \mid s).
\]

\medskip

This introduces both expressive power and new sources of instability.

\subsection{Function Approximation in Reinforcement Learning}

\begin{definition}[Value Function Approximation]
A value function is approximated by a parameterized model $V_\theta : \mathcal{S} \to \mathbb{R}$.
\end{definition}

\begin{definition}[Action-Value Approximation]
\[
Q_\theta : \mathcal{S} \times \mathcal{A} \to \mathbb{R}.
\]
\end{definition}

\begin{definition}[Policy Approximation]
A policy is parameterized as
\[
\pi_\theta(a \mid s).
\]
\end{definition}

\medskip

\noindent
\textbf{Representation role.}  
In deep RL, these functions typically factor as:
\[
s \xmapsto{\;\phi_\theta\;} z \xmapsto{\;g_\theta\;} \text{value or policy}.
\]

Thus, learning control is also learning representations.

\subsection{Value-Based Methods and Bootstrapping}

Value-based methods aim to approximate the Bellman optimality equation:
\[
Q^*(s,a) = R(s,a) + \gamma \mathbb{E}_{s'} \max_{a'} Q^*(s',a').
\]

\medskip

\noindent
\textbf{Temporal-difference (TD) learning.}  
Given a sample transition $(s,a,r,s')$, define the target:
\[
y = r + \gamma \max_{a'} Q_\theta(s',a').
\]

Then update $\theta$ to minimize:
\[
\mathcal{L}(\theta) = \bigl(Q_\theta(s,a) - y\bigr)^2.
\]

\begin{remark}[Bootstrapping]
The target depends on the current estimate $Q_\theta$.  
Learning thus uses its own predictions as targets.
\end{remark}

\paragraph{DQN intuition.}
Deep Q-Networks stabilize this process by:
\begin{itemize}
\item using a \emph{target network} for $y$
\item sampling from a replay buffer to decorrelate data
\end{itemize}

\subsection{Policy-Based Methods}

Instead of learning value functions, one may directly optimize policies.

\begin{definition}[Objective]
\[
J(\theta) = \mathbb{E}_{\pi_\theta}\left[ \sum_{t=0}^\infty \gamma^t R(s_t, a_t) \right].
\]
\end{definition}

\begin{theorem}[Policy Gradient Identity]
\[
\nabla_\theta J(\theta)
=
\mathbb{E}_{\pi_\theta}
\left[
\nabla_\theta \log \pi_\theta(a \mid s)\, Q^{\pi_\theta}(s,a)
\right].
\]
\end{theorem}

\begin{proof}[Sketch]
Write $J(\theta)$ as an expectation over trajectories.  
Differentiate under the integral and apply the log-derivative trick:
\[
\nabla_\theta \pi_\theta = \pi_\theta \nabla_\theta \log \pi_\theta.
\]
Rearranging yields the result.
\end{proof}

\begin{remark}
In practice, $Q^{\pi_\theta}$ is replaced by sampled returns or learned critics.
\end{remark}

\subsection{Actor--Critic Methods}

Actor--critic methods combine value and policy learning:
\begin{itemize}
\item \textbf{actor:} policy $\pi_\theta$
\item \textbf{critic:} value function $V_\phi$ or $Q_\phi$
\end{itemize}

\medskip

\noindent
\textbf{Interpretation.}  
The critic provides a learned representation of long-term outcomes, guiding the actor’s updates.

\subsection{Model-Free vs Model-Based Reinforcement Learning}

\begin{definition}[Model-Free RL]
Learns value functions or policies without explicitly modeling $P(s' \mid s,a)$.
\end{definition}

\begin{definition}[Model-Based RL]
Learns or uses a model of the environment:
\[
\hat{P}(s' \mid s,a), \quad \hat{R}(s,a).
\]
\end{definition}

\medskip

\noindent
\textbf{Comparison.}
\begin{itemize}
\item Model-free: simpler, often more robust, data-hungry
\item Model-based: more sample-efficient, but sensitive to model error
\end{itemize}

\begin{remark}
Modern approaches often combine both through learned latent dynamics.
\end{remark}

\subsection{Worked Example: Learned State Representation}

Consider an agent observing high-dimensional input $x_t$ (e.g., images).  
A learned encoder $\phi_\theta$ produces:
\[
z_t = \phi_\theta(x_t).
\]

\medskip

Suppose:
\[
Q_\theta(x_t, a) = g_\theta(z_t, a).
\]

If $\phi_\theta$ is trained jointly with $Q_\theta$, it learns to:
\begin{itemize}
\item discard irrelevant visual detail
\item retain task-relevant structure (e.g., object positions)
\end{itemize}

\medskip

\noindent
\textbf{Interpretation.}  
The representation $\phi_\theta$ is shaped by control objectives rather than reconstruction or prediction alone.

\subsection{Deep RL Training Loop}

A typical deep RL loop consists of:
\begin{enumerate}
\item collect transitions $(s,a,r,s')$
\item store in replay buffer
\item sample mini-batches
\item update value/policy networks
\item periodically update target networks
\end{enumerate}

\medskip

\noindent
\textbf{Interpretation.}  
Learning interleaves data collection, representation learning, and control optimization.

\subsection{Connection to Representation Learning}

Deep RL integrates representation learning with decision making:
\begin{itemize}
\item representations must support prediction of future outcomes
\item invariances are task-dependent (control-relevant)
\item transfer requires reusable state abstractions
\end{itemize}

\medskip

\noindent
\textbf{Key insight.}  
Unlike supervised learning, representations are evaluated by their effect on long-term return.

\subsection{What to Remember}

\begin{itemize}
\item Deep RL uses function approximation for value and policy
\item Bootstrapping introduces both efficiency and instability
\item Policy gradients provide a principled optimization framework
\item Representations are shaped by control objectives
\end{itemize}

\subsection{Common Confusion}

\begin{itemize}
\item Value-based and policy-based methods are not mutually exclusive
\item Bootstrapping is not the same as supervised learning
\item Good representations in RL are not necessarily interpretable
\end{itemize}

\subsection{Exercises}

\begin{exercise}
Derive the TD update for $V_\theta$ using the Bellman expectation equation.
\end{exercise}

\begin{exercise}
Show how the policy gradient identity changes when using a baseline.
\end{exercise}

\begin{exercise}
Construct an example where model-based RL fails due to model bias.
\end{exercise}

\begin{exercise}
Explain how representation learning differs between supervised learning and RL.
\end{exercise}

\subsection{Notes and References}

Key directions include:
\begin{itemize}
\item temporal-difference learning and Q-learning
\item policy gradient methods (REINFORCE)
\item deep Q-networks (DQN)
\item actor--critic methods
\item model-based RL and latent dynamics models
\end{itemize}

Deep RL remains less theoretically understood than supervised learning, particularly in the presence of function approximation and non-stationary data.